	\section{Shock Results}
\label{feartable}
\begin{tabularx}{\textwidth}{|l|X|}
	\hline
	$\le$ 3 &
		Shivering, the Subject takes a -10 penalty on all tests
		until he passes a \emph{Courage} check.
		He can spend one action to attempt this check.
		\\ \hline
	4-5   & Severely shaken, the Subject takes a -15 penalty
		on all tests for the rest of the encounter or scene
		and can't use maneuvers or transform actions.
		\\ \hline
	6-7   & Dumbstruck by fear,
			the Subject freezes in place and can’t act until he passes a Courage check at the end of one of his turns.
		This check begins at the same penalty as the initial one and gets a cumulative +10 for every previous attempt.
		He may take reactions to defend himself or run away
			(at an increased risk of stumbling as quickly avoiding obstacles is “acting”...)
		\\ \hline
	8-9 &
		Frozen by terror, the Subject is completely unable to move
		until he passes a \emph{Courage} check at the end of a following round.
		This check begins at a penalty equal to 10 times the modifier on the Shock roll
		and gets a cumulative +5 for every previous attempt.
		These Courage checks are disallowed for the first 1D5-1 rounds after fear strikes.
		He may throw himself into cover the moment fear strikes to sit there and cower.
		\\ \hline
	10-11  & Overcome by horror, the Subject loses consciousness for 1D10 minutes.
	\\ \hline
	12-13 &
		Pure horror pierces the Subject’s heart.
		His courage is reduced by 2D10+5 for a week,
		recovering by 1 each day afterwards.
		He takes a second shock effect at a value of D10.
		\\ \hline
	14-15  & Stuck in their worst nightmare, the Subject falls comatose for a day.
		\\ \hline
	16-19  & Fear breaks the Subject’s mind.
		He falls into a coma for half a week and his memory of the last 2D10 days is erased.
		After such a traumatic experience, the Subject likely develops a mental disorder (a Bane worth 2-3 GP).
		Voice and Adjudicator should discuss what this might be.
		\\ \hline
	20+ &
		The Subject suffers sudden cardiac arrest.
		He will die without immediate medical attention.
		\\ \hline
\end{tabularx}
